%==============================================================================
% PARTIDA 02: SALIDAS PARA ALUMBRADO, TOMACORRIENTES, FUERZA
% Codigos: OE.5.2.1.1.1, OE.5.2.1.1.2, OE.5.2.1.2.1, OE.5.2.1.2.2, OE.5.2.1.2.3
%==============================================================================
\subsection{OE.5.2.1.1.1 Caja Octogonal FG 100x50 mm -- Salida de Alumbrado de Techo}

\subsubsection{Especificaciones Tecnicas}

Caja octogonal de fierro galvanizado (FG) de 100 mm x 50 mm (4'' x 2''), disenada para alojar las conexiones de luminarias de techo. Debe cumplir con la NTP 370.053 (Seguridad electrica -- Cajas de derivacion metalicas) y la NTP-IEC 60670 (Cajas y envolventes para accesorios electricos).

\begin{itemize}[nosep]
  \item \textbf{Material:} Fierro galvanizado por inmersion en caliente (proteccion anticorrosiva Z275 segun NTP-ISO 1461).
  \item \textbf{Especificaciones dimensionales:} 100 mm ± 2 mm entre caras opuestas, profundidad 50 mm ± 1 mm.
  \item \textbf{Orejas de fijacion:} Dos orejas roscadas (1/4''-20 UNC) para soporte de luminaria, separacion 75 mm ± 1 mm.
  \item \textbf{Entradas de tuberia:} 4 entradas ciegas (knockouts) de 3/4'' en caras laterales, mas 1 entrada central en la base para alimentacion.
  \item \textbf{Carga maxima admisible:} 25 kg en orejas de fijacion.
  \item \textbf{Grado de proteccion:} IP-30 instalada con tapa ciega.
\end{itemize}

\paragraph{Normativa aplicable:}
\begin{itemize}[nosep]
  \item CNE-U Articulo 240.1: Las cajas de derivacion deben ser accesibles y de material aprobado.
  \item CNE-U Articulo 240.3: Volumen interior minimo de cajas segun cantidad de conductores.
  \item RNE EM.010 Articulo 12: Ubicacion y fijacion de cajas de salida.
\end{itemize}

\subsubsection{Requisitos de Materiales}

\begin{longtable}{L{5.0cm} L{2.0cm} L{7.5cm}}
\toprule
\textbf{Material} & \textbf{Cant.} & \textbf{Especificacion tecnica} \\
\midrule
Caja octogonal FG 100x50 mm & 19 pza & Fierro galvanizado, 4 entradas 3/4'', con 2 orejas roscadas \\
Tapa ciega metalica 4'' & 19 pza & FG con tornillo de fijacion 1/4''-20 UNC \\
Conector tipo "romex" 3/4'' & 38 pza | Para ingreso de tuberia PVC, cuerpo metalico con contratuerca \\
Tornillo de fijacion 1/4'' x 1'' & 76 pza | Para anclaje de orejas de luminaria (2 por caja) \\
Cinta aislante vinilica 3M & 2 pza | Para proteccion temporal de conexiones \\
\bottomrule
\end{longtable}

\subsubsection{Metodo de Ejecucion}

\begin{enumerate}
  \item \textbf{Replanteo:} Marcar en la losa de techo la ubicacion exacta de cada caja octogonal segun planos IE-02, IE-03 e IE-04. La posicion debe coincidir con el centro geometrico del ambiente o con la ubicacion disenada de la luminaria. Tolerancia: ±5 cm.
  \item \textbf{Preparacion de la caja:} Retirar los knockouts de las entradas necesarias (2 entradas: una para tuberia de alimentacion, otra para tuberia de derivacion si corresponde). Limpiar rebabas con lima.
  \item \textbf{Fijacion a losa:}
  \begin{enumerate}
    \item Colocar la caja octogonal en la posicion marcada.
    \item Fijar a la losa de techo mediante alambre calibre 16 amarrado a la armadura de acero, o mediante clavos de 2'' en encofrado de madera.
    \item La cara frontal de la caja debe quedar al ras con el fondo de la losa (para losa maciza) o al ras con el falso cielo raso.
    \item Nivelar la caja en ambos ejes (tolerancia: ±2 mm).
  \end{enumerate}
  \item \textbf{Conexion de tuberias:}
  \begin{enumerate}
    \item Insertar los conectores tipo "romex" en las entradas de la caja.
    \item Fijar con contratuerca del lado interior de la caja, apretando con pinza.
    \item Insertar la tuberia PVC de 3/4'' en el conector hasta el tope.
    \item Sellar la union con pegamento PVC.
  \end{enumerate}
  \item \textbf{Dejado de conductores:}
  \begin{enumerate}
    \item Pasar los conductores (fase, neutro y tierra) desde la caja de derivacion hasta la caja octogonal.
    \item Dejar 0.30 m de conductor desde el borde de la caja para conexiones futuras de la luminaria.
    \item Doblar los conductores dentro de la caja sin forzar las curvas.
  \end{enumerate}
  \item \textbf{Colocacion de tapa ciega:}
  \begin{enumerate}
    \item Instalar la tapa ciega metalica sobre la caja, fijando con tornillo 1/4''-20 UNC a la oreja central.
    \item La tapa permanecera durante toda la etapa de construccion para proteger la caja del ingreso de concreto, polvo y humedad.
  \end{enumerate}
  \item \textbf{Identificacion:} Marcar temporalmente el circuito correspondiente (C1, C3 o C5) en la caja con cinta de color o marcador indeleble.
\end{enumerate}

\subsubsection{Control de Calidad}

\begin{longtable}{L{4.5cm} L{4.5cm} L{4.5cm}}
\toprule
\textbf{Parametro} & \textbf{Metodo de verificacion} & \textbf{Criterio de aceptacion} \\
\midrule
Posicion & Medicion con cinta metrica desde referencias fijas & ±5 cm respecto al plano \\
Nivelacion & Nivel de burbuja de 30 cm sobre la caja & ±2 mm en ambos ejes \\
Fijacion & Prueba de traccion manual (5 kg) & Sin desplazamiento \\
Altura de conductores & Medicion con cinta metrica desde borde de caja & 0.30 m ± 2 cm \\
Conexion de tuberias & Inspeccion visual y prueba de traccion (2 kg) & Tuberia firmemente insertada \\
\bottomrule
\end{longtable}

\subsubsection{Metodo de Medicion}

Se medira por pieza (pza) de caja octogonal de alumbrado de techo completamente instalada, incluyendo caja, tapa ciega, conectores, fijacion a losa, conexion de tuberias, dejado de conductores y proteccion temporal.

\subsubsection{Unidad de Medida}

pza (pieza).

\subsubsection{Forma de Pago}

Pago por pieza instalada, verificada en replanteo (posicion y nivelacion) y aprobada por supervision antes del vaciado de losa o cierre de cielo raso.

%------------------------------------------------------------------------------
\subsection{OE.5.2.1.1.2 Caja Rectangular FG 4''x2'' -- Salida para Interruptores}

\subsubsection{Especificaciones Tecnicas}

Caja rectangular de fierro galvanizado (FG) de 4'' x 2'' (100 mm x 50 mm) para instalacion de interruptores de luz unipolar (S1) o de conmutacion (S3). La caja debe cumplir con la NTP 370.053 y NTP-IEC 60670.

\begin{itemize}[nosep]
  \item \textbf{Material:} Fierro galvanizado, espesor nominal 1.2 mm.
  \item \textbf{Dimensiones:} 100 mm x 50 mm x 50 mm (largo x ancho x profundidad).
  \item \textbf{Orejas:} Dos orejas roscadas 1/4''-20 UNC para fijacion del mecanismo del interruptor.
  \item \textbf{Entradas:} 2 knockouts de 3/4'' en extremos, 2 knockouts en parte posterior.
  \item \textbf{Altura de instalacion:} 1.20 m sobre el NPT (Nivel de Piso Terminado), medido desde el borde superior de la caja.
  \item \textbf{Capacidad maxima de conductores:} 4 conductores de 2.5 mm2 segun CNE-U Tabla 240.1.
\end{itemize}

\subsubsection{Requisitos de Materiales}

\begin{longtable}{L{5.0cm} L{2.0cm} L{7.5cm}}
\toprule
\textbf{Material} & \textbf{Cant.} & \textbf{Especificacion tecnica} \\
\midrule
Caja rectangular FG 4''x2'' & 10 pza & FG galvanizado, 2 orejas roscadas, 4 knockouts 3/4'' \\
Conector tipo "romex" 3/4'' & 20 pza & Metalico con contratuerca, para tuberia PVC \\
Tornillo de anclaje 1/4'' x 2'' con tarugo & 20 pza | Para fijacion a muro de ladrillo \\
Tapa ciega provisional & 10 pza | Plastica o metalica para proteccion durante obra \\
Nivel de burbuja tipo torpedo & 1 pza | 30 cm, para verificacion de instalacion \\
\bottomrule
\end{longtable}

\subsubsection{Metodo de Ejecucion}

\begin{enumerate}
  \item \textbf{Replanteo:} Marcar en el muro la ubicacion del interruptor segun plano IE-01. Verificar:
  \begin{itemize}
    \item Altura: 1.20 m sobre NPT (medido desde borde superior de caja instalada).
    \item Distancia minima al marco de puerta: 0.15 m del lado del tirador.
    \item No debe quedar detras de puertas, muebles fijos o artefactos.
  \end{itemize}
  \item \textbf{Apertura de caja en muro:}
  \begin{enumerate}
    \item Realizar la excavacion en muro de ladrillo con comba y cincel o cortadora electrica con disco de corte.
    \item Dimensiones de la excavacion: 110 mm x 60 mm x 55 mm (profundidad).
    \item Limpiar el polvo y residuos con brocha.
  \end{enumerate}
  \item \textbf{Instalacion de caja:}
  \begin{enumerate}
    \item Insertar la caja rectangular en el hueco.
    \item La cara frontal debe quedar 2 mm por debajo de la superficie del tarrajeo terminado (para permitir que la placa frontal quede al ras).
    \item Fijar con tornillo 1/4'' x 2'' y tarugo de expansion en las esquinas.
    \item Verificar nivelacion vertical y horizontal con nivel de burbuja (tolerancia: ±1 mm).
  \end{enumerate}
  \item \textbf{Conexion de tuberia:}
  \begin{enumerate}
    \item Retirar knockouts de los extremos necesarios.
    \item Instalar conectores "romex" y fijar con contratuerca.
    \item Insertar tuberia PVC de 3/4'' y sellar con pegamento.
  \end{enumerate}
  \item \textbf{Dejado de conductores:}
  \begin{enumerate}
    \item Pasar conductores desde caja de derivacion hasta caja de interruptor.
    \item Para interruptor S1 (unipolar): conductor de fase y conductor de retorno.
    \item Para interruptor S3 (conmutacion): fase, dos conductores viajeros y retorno.
    \item Dejar 0.25 m de conductor desde el borde de la caja.
  \end{enumerate}
  \item \textbf{Proteccion:} Colocar tapa ciega provisional para evitar ingreso de mortero durante el tarrajeo.
\end{enumerate}

\subsubsection{Control de Calidad}

\begin{longtable}{L{4.5cm} L{4.5cm} L{4.5cm}}
\toprule
\textbf{Parametro} & \textbf{Metodo de verificacion} & \textbf{Criterio de aceptacion} \\
\midrule
Altura de instalacion & Cinta metrica desde NPT hasta borde superior caja & 1.20 m ± 1 cm \\
Distancia a marco de puerta & Cinta metrica & Minimo 0.15 m \\
Nivelacion vertical & Nivel de burbuja sobre cara frontal & ±1 mm \\
Nivelacion horizontal & Nivel de burbuja sobre borde superior & ±1 mm \\
Empotramiento & Regla de aluminio de 1 m sobre la caja & 2 mm ± 1 mm bajo superficie \\
Fijacion & Prueba de traccion manual (3 kg) & Sin desplazamiento \\
\bottomrule
\end{longtable}

\subsubsection{Metodo de Medicion}

Se medira por pieza (pza) de caja rectangular para interruptor completamente instalada, incluyendo excavacion, fijacion, conexion de tuberia y dejado de conductores.

\subsubsection{Unidad de Medida}

pza (pieza).

\subsubsection{Forma de Pago}

Pago por pieza instalada y verificada antes del tarrajeo definitivo.

%------------------------------------------------------------------------------
\subsection{OE.5.2.1.2.1 Caja Rectangular FG 4''x2'' -- Salida de Tomacorriente Doble con Tierra}

\subsubsection{Especificaciones Tecnicas}

Caja rectangular FG 4''x2'' para alojar tomacorriente doble bipolar con contacto a tierra, 15 A -- 250 V, tipo universal. Debe cumplir con NTP 370.053, NTP-IEC 60884-1 (Clavijas y tomacorrientes) y NTP 370.304.

\begin{itemize}[nosep]
  \item \textbf{Altura de instalacion:}
  \begin{itemize}
    \item Ambientes generales: 0.40 m sobre NPT (medido desde borde superior de caja).
    \item Cocina (sobre encimera): 1.10 m sobre NPT.
    \item Lavanderia: 1.30 m sobre NPT.
  \end{itemize}
  \item \textbf{Conductores requeridos:} 3 conductores de 2.5 mm2 LSOH: fase (rojo/marron), neutro (blanco/celeste), proteccion a tierra (verde/verde-amarillo).
  \item \textbf{Capacidad del tomacorriente:} 15 A, 250 V, con borne de tierra tipo pin de 3/8''.
\end{itemize}

\subsubsection{Requisitos de Materiales}

\begin{longtable}{L{5.0cm} L{2.0cm} L{7.5cm}}
\toprule
\textbf{Material} & \textbf{Cant.} & \textbf{Especificacion tecnica} \\
\midrule
Caja rectangular FG 4''x2'' & 24 pza & FG galvanizado, 2 orejas roscadas \\
Conector "romex" 3/4'' & 48 pza | Metalico con contratuerca \\
Tornillo anclaje 1/4'' x 2'' con tarugo & 48 pza | Para fijacion a muro \\
Tomacorriente doble con tierra 15A-250V & 24 pza | TICINO o BTICINO, color blanco, con placa frontal 2 modulos \\
Placa frontal decorativa 2 modulos & 24 pza | Color blanco, ABS ignifugo, material autoextinguible \\
Conductor LSOH 2.5 mm2 rojo & Segun metrado | Fase, color rojo o marron \\
Conductor LSOH 2.5 mm2 blanco & Segun metrado | Neutro, color blanco o celeste \\
Conductor LSOH 2.5 mm2 verde & Segun metrado | Tierra, color verde o verde-amarillo \\
\bottomrule
\end{longtable}

\subsubsection{Metodo de Ejecucion}

\begin{enumerate}
  \item \textbf{Replanteo:} Marcar ubicacion segun plano IE-01 a las alturas especificadas. Verificar que no existan obstaculos (tuberias de agua, desague, gas).
  \item \textbf{Instalacion de caja:} Seguir el mismo procedimiento de la partida OE.5.2.1.1.2 (excavacion, fijacion, nivelacion).
  \item \textbf{Canalizacion:} Conectar tuberia PVC pesada (PV-P) de 3/4'' desde la caja de derivacion hasta la caja del tomacorriente.
  \item \textbf{Tendido de conductores:} Pasar los 3 conductores (fase, neutro, tierra) de 2.5 mm2. Dejar 0.25 m desde el borde de la caja.
  \item \textbf{Instalacion del mecanismo:}
  \begin{enumerate}
    \item Pelar los conductores 1.5 cm con pelacables automatico.
    \item Conectar fase (rojo) al borne L o marron del tomacorriente.
    \item Conectar neutro (blanco) al borne N o azul del tomacorriente.
    \item Conectar tierra (verde) al borne de tierra (simbolo $\perp$ o tornillo central).
    \item Apretar tornillos de bornes a 2.5 Nm con destornillador de precision.
    \item Fijar el mecanismo a la caja mediante los tornillos de las orejas.
  \end{enumerate}
  \item \textbf{Colocacion de placa frontal:} Encajar la placa frontal decorativa sobre el mecanismo y fijar con tornillos visibles (si aplica).
\end{enumerate}

\subsubsection{Control de Calidad}

\begin{longtable}{L{4.5cm} L{4.5cm} L{4.5cm}}
\toprule
\textbf{Parametro} & \textbf{Metodo de verificacion} & \textbf{Criterio de aceptacion} \\
\midrule
Polaridad & Comprobador de tomacorrientes (tester) & Fase en borne L, neutro en N, tierra conectada \\
Tension & Multimetro digital entre fase y neutro & 220 V ± 10\% (198 V -- 242 V) \\
Continuidad de tierra & Multimetro entre borne de tierra y barra de tierra del tablero & < 1 $\Omega$ \\
Corriente de fuga & Pinza amperimetrica de fuga & < 5 mA \\
Fijacion del mecanismo & Inspeccion visual y prueba de traccion del tomacorriente (2 kg) & Sin juego ni desplazamiento \\
\bottomrule
\end{longtable}

\subsubsection{Metodo de Medicion}

Se medira por pieza (pza) de tomacorriente doble con tierra completamente instalado, incluyendo caja, conductores, mecanismo, placa frontal y pruebas electricas.

\subsubsection{Unidad de Medida}

pza (pieza).

\subsubsection{Forma de Pago}

Pago por pieza instalada, verificada con comprobador de tomacorrientes (tester) y aprobada por supervision.

%------------------------------------------------------------------------------
\subsection{OE.5.2.1.2.2 Caja Rectangular FG 4''x2'' -- Salida de Tomacorriente Doble con Proteccion GFCI}

\subsubsection{Especificaciones Tecnicas}

Tomacorriente doble con interruptor diferencial incorporado (GFCI -- Ground Fault Circuit Interrupter), 15 A -- 250 V, sensibilidad 30 mA, para instalacion en ambientes humedos (banos, cocinas, lavanderias). Debe cumplir con NTP-IEC 61008-1 (Interruptores diferenciales) y NTP 370.304.

\begin{itemize}[nosep]
  \item \textbf{Corriente nominal:} 15 A.
  \item \textbf{Tension nominal:} 250 V.
  \item \textbf{Corriente diferencial residual nominal ($I_{\Delta n}$):} 30 mA.
  \item \textbf{Tiempo de disparo:} < 40 ms a 1x $I_{\Delta n}$.
  \item \textbf{Proteccion:} Tipo A (protege corrientes alternas y pulsantes).
  \item \textbf{Altura de instalacion:} 1.20 m sobre NPT en banos (fuera de la zona de salpicadura directa, minimo 0.60 m de la ducha).
  \item \textbf{Norma:} NTP-IEC 61008-1, UL 943.
\end{itemize}

\paragraph{Ambientes donde se instala GFCI (segun CNE-U Articulo 560.2):}
\begin{itemize}[nosep]
  \item Banos (todos los tomacorrientes).
  \item Cocinas (tomacorrientes a menos de 1.80 m de la pileta).
  \item Lavanderias.
  \item Exteriores.
\end{itemize}

\subsubsection{Requisitos de Materiales}

\begin{longtable}{L{5.0cm} L{2.0cm} L{7.5cm}}
\toprule
\textbf{Material} & \textbf{Cant.} & \textbf{Especificacion tecnica} \\
\midrule
Caja rectangular FG 4''x2'' profunda & 2 pza | Profundidad minima 60 mm para alojar mecanismo GFCI \\
Conector "romex" 3/4'' & 4 pza | Metalico con contratuerca \\
Tomacorriente GFCI 15A-250V-30mA & 2 pza | TICINO, LEGRAND o similar, con bornes LINE/LOAD \\
Placa frontal decorativa 2 modulos & 2 pza | Color blanco, ABS ignifugo \\
Conductor LSOH 2.5 mm2 rojo & Segun metrado | Fase \\
Conductor LSOH 2.5 mm2 blanco & Segun metrado | Neutro \\
Conductor LSOH 2.5 mm2 verde & Segun metrado | Tierra \\
\bottomrule
\end{longtable}

\subsubsection{Metodo de Ejecucion}

\begin{enumerate}
  \item \textbf{Replanteo:} Ubicar en el muro del bano a 1.20 m sobre NPT, minimo 0.60 m de la ducha o banera.
  \item \textbf{Instalacion de caja profunda:} Seguir procedimiento del tomacorriente general, utilizando caja de profundidad minima 60 mm.
  \item \textbf{Cableado:} Pasar conductores de 2.5 mm2 (fase, neutro, tierra) desde la caja de derivacion.
  \item \textbf{Conexion del GFCI:}
  \begin{enumerate}
    \item Identificar los bornes del GFCI (marcados en la parte posterior):
    \begin{itemize}
      \item "LINE" o "LINE L": Borne de entrada de fase.
      \item "LINE N": Borne de entrada de neutro.
      \item "LOAD" o "LOAD L": Borne de salida de fase (solo si se protegen tomacorrientes aguas abajo).
      \item "LOAD N": Borne de salida de neutro (solo si se protegen tomacorrientes aguas abajo).
      \item Borne de tierra (verde): Conexion de proteccion a tierra.
    \end{itemize}
    \item Para este proyecto (proteccion de un solo punto): conectar fase y neutro de entrada en bornes LINE. NO conectar bornes LOAD.
    \item Conectar tierra al borne de tierra.
    \item Apretar bornes a 2.5 Nm.
  \end{enumerate}
  \item \textbf{Prueba de funcionamiento:}
  \begin{enumerate}
    \item Presionar boton "RESET" para habilitar el tomacorriente.
    \item Presionar boton "TEST": debe audirse un "clic" y el tomacorriente debe cortar el suministro.
    \item Presionar "RESET" nuevamente para restablecer.
    \item Verificar que el indicador luminoso (si lo tiene) encienda en estado normal.
  \end{enumerate}
\end{enumerate}

\subsubsection{Control de Calidad}

\begin{longtable}{L{4.5cm} L{4.5cm} L{4.5cm}}
\toprule
\textbf{Parametro} & \textbf{Metodo de verificacion} & \textbf{Criterio de aceptacion} \\
\midrule
Prueba TEST & Presionar boton TEST & Disparo inmediato (< 40 ms) \\
Prueba RESET & Presionar boton RESET & Restablecimiento correcto \\
Tiempo de disparo & Medidor de tiempo de disparo diferencial (RCD tester) & < 40 ms a 30 mA \\
Corriente de fuga & Pinza amperimetrica de fuga & < 5 mA en estado normal \\
Polaridad & Comprobador de tomacorrientes GFCI & Correcta (fase/neutro/tierra) \\
Tension & Multimetro digital & 220 V ± 10\% \\
\bottomrule
\end{longtable}

\subsubsection{Metodo de Medicion}

Se medira por pieza (pza) de tomacorriente GFCI instalado, incluyendo caja profunda, conductores, mecanismo, placa y pruebas de funcionamiento (TEST/RESET).

\subsubsection{Unidad de Medida}

pza (pieza).

\subsubsection{Forma de Pago}

Pago por pieza instalada y verificada mediante prueba de disparo (TEST) y restablecimiento (RESET), con protocolo de prueba firmado.

%------------------------------------------------------------------------------
\subsection{OE.5.2.1.2.3 Caja Rectangular FG 4''x2'' -- Salida de Tomacorriente Servicio Pesado 20A}

\subsubsection{Especificaciones Tecnicas}

Tomacorriente doble de servicio pesado 20 A -- 250 V, con contacto a tierra de mayor diametro (espiga de 5/16''), para cargas intensivas como electrobomba, lavadora, horno microondas o equipos de cocina de alta potencia. Debe cumplir con NTP-IEC 60884-1.

\begin{itemize}[nosep]
  \item \textbf{Corriente nominal:} 20 A.
  \item \textbf{Tension nominal:} 250 V.
  \item \textbf{Configuracion:} 2 polos + tierra (NEMA 5-20R).
  \item \textbf{Conductor requerido:} 2.5 mm2 LSOH minimo (segun CNE-U Tabla 5.2 para 20 A).
  \item \textbf{Proteccion termomagnetica:} 2P-20A (curva C).
\end{itemize}

\subsubsection{Requisitos de Materiales}

\begin{longtable}{L{5.0cm} L{2.0cm} L{7.5cm}}
\toprule
\textbf{Material} & \textbf{Cant.} & \textbf{Especificacion tecnica} \\
\midrule
Caja rectangular FG 4''x2'' & 3 pza | FG galvanizado, 2 orejas roscadas \\
Tomacorriente servicio pesado 20A-250V & 3 pza | Doble, con tierra, configuracion NEMA 5-20R \\
Placa frontal decorativa 2 modulos & 3 pza | ABS ignifugo, color blanco \\
Conductor LSOH 2.5 mm2 (fase+neutro+tierra) & Segun metrado | Para alimentacion dedicada \\
\bottomrule
\end{longtable}

\subsubsection{Metodo de Ejecucion}

Identico al procedimiento del tomacorriente general (OE.5.2.1.2.1), con las siguientes particularidades:
\begin{itemize}
  \item Verificar que la proteccion termomagnetica del circuito sea 2P-20A.
  \item Verificar que el tomacorriente tenga grabado "20A" en el cuerpo.
  \item Conexion de conductores de 2.5 mm2 con terminales de horquilla para mejor contacto.
\end{itemize}

\subsubsection{Control de Calidad}

\begin{itemize}
  \item Verificar que el tomacorriente soporte 20 A nominales (inspeccion del sello o grabado en el cuerpo).
  \item Comprobar que la espiga de tierra del tomacorriente sea de 5/16'' (mayor que la estandar de 3/16'').
  \item Verificar que la proteccion termomagnetica asociada sea de 20 A (curva C).
  \item Pruebas de polaridad, tension y continuidad de tierra identicas al tomacorriente general.
\end{itemize}

\subsubsection{Metodo de Medicion}

Se medira por pieza (pza) de tomacorriente servicio pesado 20 A completamente instalado.

\subsubsection{Unidad de Medida}

pza (pieza).

\subsubsection{Forma de Pago}

Pago por pieza instalada y verificada, incluyendo comprobacion de la proteccion termomagnetica asociada.
